\Needspace{6\baselineskip}
\begin{center}\hypertarget{native-266-article192}{}\pdfbookmark[2]{192.}{article-192}192.\end{center}
Das im vorigen Artikel dargelegte Verfahren eignet sich vor allem für die
\emph{erste} Bestimmung einer Bahn unter Berücksichtigung der Störungen.
Sobald sowohl die mittleren elliptischen Elemente als auch die
Gleichungen für die Störungen bereits näherungsweise bekannt sind, lässt
sich die genaueste Bestimmung mithilfe möglichst vieler Beobachtungen am
bequemsten nach dem Verfahren aus Art. 187 durchführen, das hier keiner
besonderen Erläuterung bedarf. Liegen hinreichend viele Beobachtungen
von höchster Güte vor und umfassen sie einen langen Zeitraum, so kann
dieses Verfahren in vielen Fällen zugleich zur genaueren Bestimmung
der Massen der störenden Planeten dienen, zumindest der größeren unter
ihnen. Erscheint nämlich die bei der Berechnung der Störungen
angenommene Masse eines störenden Planeten noch nicht hinreichend
gesichert, so führt man neben den sechs von den Korrekturen der Elemente
abhängigen Unbekannten eine weitere Unbekannte $\mu$ ein. Dazu setzt man
das Verhältnis der korrigierten zur angenommenen Masse gleich $1+\mu$
zu $1$. Man darf dann annehmen, dass sich die Störungen selbst in
demselben Verhältnis ändern. Dadurch entsteht offenbar in jeder
berechneten Position ein neuer linearer Term, der $\mu$ enthält und
sich ohne Schwierigkeit entwickeln lässt. Der Vergleich der berechneten
mit den beobachteten Positionen nach den oben dargelegten
\SourceLeaf{267}Grundsätzen liefert zusammen mit den Korrekturen der Elemente
auch die Korrektur $\mu$. Auf diese Weise kann man sogar die Massen
\emph{mehrerer} Planeten genauer bestimmen, sofern sie hinreichend
beträchtliche Störungen ausüben. Zweifellos werden die Bewegungen der
neuen Planeten, vor allem von Pallas und Juno, die durch Jupiter so
starke Störungen erfahren, nach einigen Jahrzehnten auf diese Weise
eine äußerst genaue Bestimmung der Masse Jupiters ermöglichen. Vielleicht
lässt sich sogar die Masse des einen oder anderen dieser neuen
Planeten selbst eines Tages aus den Störungen bestimmen, die er auf die
übrigen ausübt.
\begin{center}\rule{3cm}{.4pt}\end{center}
